\documentclass[11pt]{article}
\usepackage[margin=1in]{geometry}
\usepackage{amsmath,amssymb}
\usepackage{booktabs}
\usepackage{hyperref}

\title{Step-by-Step Monte Carlo Procedure\\
\large \texttt{MC\_immune\_response.py}}
\author{}
\date{}

\begin{document}
\maketitle

\noindent This document walks through the Monte Carlo (MC) procedure in
\texttt{MC\_immune\_response.py} in the exact order it executes, referencing
the variable and function names used in the code so it can be read alongside
the file.

\begin{enumerate}

\item \textbf{Load the model.} \texttt{model\_Bacteremia.py} is loaded
dynamically via \texttt{importlib} into \texttt{model\_mod}, giving access
to \texttt{PharmacokineticModel}, \texttt{ImmuneResponse}, and
\texttt{dual\_reservoir\_model}.

\item \textbf{Fix run-level constants.} \texttt{NUM\_ITER = 1000},
\texttt{LOD = 10.0} CFU/mL, \texttt{TOTAL\_H = 1944} h,
\texttt{VANCO\_START = 504} h, and a shared evaluation grid
\texttt{t\_eval = linspace(0, 1944, 1450)}, reused by every ODE solve in the
script.

\item \textbf{Build the drug-concentration callables once.}
\texttt{PharmacokineticModel().concentration\_function(...)} is called once
for vancomycin and once for linezolid, producing \texttt{van\_func(t)} and
\texttt{lzd\_func(t)}. These are reused unmodified in every solve below.

\item \textbf{Fix all parameters except $k_{\mathrm{immune}}$.} A single
\texttt{params} dictionary is built once (values in Table~\ref{tab:params})
and used, unmodified, by every solve in the script. The initial state
\texttt{Y0 = [0.0, 0.0, 100.0, 100.0]} ($S_b, R_b, S_{res}, R_{res}$) is
likewise fixed for every iteration. $k_{\mathrm{immune}}$ is deliberately
absent from \texttt{params}: it lives instead in a separate
\texttt{ImmuneResponse} object that is reconstructed per solve (steps 8 and
11), since the immune-clearance rate is stored on that object rather than
in the parameter dictionary consumed by \texttt{dual\_reservoir\_model}.

\begin{table}[h]
\centering
\begin{tabular}{@{}ll@{}}
\toprule
Parameter & Value \\
\midrule
$\rho_S$ & $0.61$ h$^{-1}$ \\
$\rho_R$ & $0.55$ h$^{-1}$ \\
$\rho_{res,S}$ & $0.175$ h$^{-1}$ \\
$\rho_{res,R}$ & $0.1765$ h$^{-1}$ \\
$E_{max,V}$ & $0.40$ \\
$EC_{50,V}$ & $0.245$ \\
$E_{max,L}$ & $0.8$ \\
$EC_{50,L}$ & $1.0$ \\
$B_{max,blood}$ & $6000$ CFU/mL \\
$B_{max,res}$ & $10^4$ CFU/mL \\
$\mathtt{van\_res\_fraction}$ & $0.15$ \\
$\mathtt{lzd\_res\_fraction}$ & $0.45$ \\
$f_{r\to b}$ & $5\times10^{-5}$ h$^{-1}$ \\
$f_{b\to r}$ & $1\times10^{-5}$ h$^{-1}$ \\
\bottomrule
\end{tabular}
\caption{The fixed \texttt{params} dictionary. Every value here is identical across all iterations and both the switch-curve sweep and the Monte Carlo loop.}
\label{tab:params}
\end{table}

\item \textbf{Compute the log-normal sampling parameters.} With
\texttt{K\_IMMUNE\_MEAN = 0.172} and \texttt{CV = 0.25}:
\begin{equation}
\sigma_{\ln} = \sqrt{\ln(1+CV^2)}, \qquad
\mu_{\ln} = \ln(\texttt{K\_IMMUNE\_MEAN}) - \tfrac{1}{2}\sigma_{\ln}^2.
\end{equation}

\item \textbf{Draw all 1000 samples at once.} \texttt{np.random.seed(42)}
seeds NumPy's legacy global random state, then
\texttt{k\_immune\_samples = np.random.lognormal(}$\mu_{\ln}, \sigma_{\ln},$
\texttt{NUM\_ITER)} draws the full length-1000 vector in a single call.
Sample $i$ of this vector is used by Monte Carlo iteration $i$ in step 11.
The realized mean, standard deviation, $CV$, and 5th/95th percentiles of the
drawn sample are printed immediately.

\item \textbf{Compute the deterministic switch curve (101 points).} Before
the Monte Carlo loop runs, a helper \texttt{\_final\_counts(k\_immune)}
constructs a fresh \texttt{ImmuneResponse(k\_immune=k\_immune)}, integrates
\texttt{dual\_reservoir\_model} with \texttt{odeint}
(\texttt{rtol=1e-7}, \texttt{atol=1e-9}, \texttt{mxstep=5000}) from the fixed
\texttt{Y0} over the fixed \texttt{t\_eval}, and returns only the final-time
$R_b$ and $R_{res}$ (\texttt{sol[-1, 1]}, \texttt{sol[-1, 3]}). This is
evaluated at \texttt{sweep\_k\_immune = linspace(0.02, 0.22, 101)}, storing
the results in \texttt{sweep\_R\_b} and \texttt{sweep\_R\_res}.

\item \textbf{Root-find \texttt{ESCAPE\_THRESH}.} Using the same
\texttt{\_final\_counts}, \texttt{scipy.optimize.brentq} solves
$R_{res}^{\mathrm{final}}(k_{\mathrm{immune}}) - \mathrm{LOD} = 0$ over
$[0.02, 0.22]$ (falling back to the coarse sweep crossing if the endpoint
signs don't bracket a root), giving a single precise threshold value,
\texttt{ESCAPE\_THRESH}.

\item \textbf{Pre-allocate the history arrays.} Four arrays
\texttt{S\_b\_hist}, \texttt{R\_b\_hist}, \texttt{S\_res\_hist},
\texttt{R\_res\_hist}, each of shape $1000 \times 1450$, are filled with
\texttt{NaN} as placeholders to be overwritten row-by-row in step 11.

\item \textbf{Main Monte Carlo loop (1000 iterations).} For each sampled
value \texttt{k\_im} in \texttt{k\_immune\_samples} (iteration index $i$):
\begin{enumerate}
\item[(a)] build \texttt{immune\_i = ImmuneResponse(k\_immune=k\_im)} -- a
fresh immune-response object carrying this iteration's sampled rate (the
shared \texttt{params} dictionary itself is not touched);
\item[(b)] integrate \texttt{dual\_reservoir\_model} from the fixed
\texttt{Y0} over the fixed \texttt{t\_eval}, with \texttt{odeint}
(\texttt{rtol=1e-7}, \texttt{atol=1e-9}, \texttt{mxstep=5000}), passing
\texttt{immune\_i} in place of a shared immune model;
\item[(c)] clip all four state trajectories to be non-negative
(\texttt{np.clip(sol[:, j], 0, None)});
\item[(d)] replace every value below \texttt{LOD} with \texttt{NaN}
(\texttt{np.where(x < LOD, np.nan, x)}) for each of the four compartments --
this is this script's sentinel for "below detection," used later so that
\texttt{np.nanpercentile} and the log-scale plots simply skip those points;
\item[(e)] store the resulting length-1450 trajectory for each compartment
into row $i$ of the four history arrays;
\item[(f)] print progress every 50 completed iterations.
\end{enumerate}
No parameter other than $k_{\mathrm{immune}}$ differs between iterations.

\item \textbf{Render the 4-panel kinetics figure.} For each of the four
compartments ($S_b$, $R_b$, $S_{res}$, $R_{res}$), the helper
\texttt{\_draw\_panel}:
\begin{enumerate}
\item[(a)] shades the vancomycin and linezolid treatment windows and draws
the LOD reference line;
\item[(b)] selects \texttt{N\_TRACES = 40} of the 1000 stored trajectories
at random, using a \emph{separately seeded} generator
(\texttt{np.random.default\_rng(0).choice(NUM\_ITER, 40, replace=False)},
independent of the \texttt{np.random.seed(42)} state used for sampling
$k_{\mathrm{immune}}$), and plots each as a faint individual line;
\item[(c)] computes the 5th/50th/95th percentile across all 1000 rows at
each time point via \texttt{np.nanpercentile} (which ignores the
\texttt{NaN}-masked sub-LOD entries from step 11d) and plots the
5th--95th band and the median line.
\end{enumerate}
The resulting figure is saved as \texttt{mc\_immune\_response.png}.

\item \textbf{Compute per-iteration final outcomes.}
\texttt{final\_R\_b} and \texttt{final\_R\_res} take the last time point of
each row of \texttt{R\_b\_hist}/\texttt{R\_res\_hist}
(\texttt{hist[:, -1]}), mapping any remaining \texttt{NaN} back to $0$. A
boolean mask \texttt{relapsed = (final\_R\_b > LOD) | (final\_R\_res > LOD)}
marks each of the 1000 iterations as relapse (resistant infection persists
in blood and/or reservoir) or resolution (both cleared); the relapse and
resolution rates are printed as counts and percentages.

\item \textbf{Render the switch-curve and population figure.} A second
figure is built in two stacked panels sharing the $x$-axis
($k_{\mathrm{immune}}$):
\begin{enumerate}
\item[(a)] the \emph{top} panel plots the deterministic
\texttt{sweep\_k\_immune} vs.\ \texttt{sweep\_R\_b}/\texttt{sweep\_R\_res}
curves from step 7 (values at or below \texttt{LOD} floored to a display
value for the log axis), with the region below \texttt{ESCAPE\_THRESH}
shaded as "escape" and the region above shaded as "suppressed" -- for
$k_{\mathrm{immune}}$, \emph{low} values are the escape zone (weak host
immunity), the reverse of the convention used for $EC_{50,L}$ in
\texttt{MC\_ec50\_lzd.py};
\item[(b)] the \emph{bottom} panel is a histogram of the same 1000 sampled
\texttt{k\_immune\_samples} from step 6, split into two subsets by each
sample's own simulated outcome from step 12
(\texttt{escaped\_mask = final\_R\_res > LOD}) -- not merely by which side
of \texttt{ESCAPE\_THRESH} the sampled value falls on -- and colored
accordingly, with the fraction of escaped vs.\ suppressed patients reported
in the legend.
\end{enumerate}
The resulting figure is saved as \texttt{mc\_immune\_response\_switch.png}.

\end{enumerate}

\section{Design rationale}

The step-by-step procedure above describes \emph{what} the script does. This
section explains \emph{why} its parameters and design choices take the
values they do, drawn from the script's module docstring and inline
comments.

\subsection{Why $k_{\mathrm{immune}}$, and why this direction}

This script isolates uncertainty in the host's innate immune clearance rate,
$k_{\mathrm{immune}}$, mirroring \texttt{MC\_ec50\_lzd.py} and
\texttt{MC\_emax\_lzd.py}'s pattern of varying one PD/host parameter while
holding everything else at baseline -- but for host response rather than
drug potency. As with $E_{max,L}$ (and unlike $EC_{50,L}$), \emph{low}
values of $k_{\mathrm{immune}}$ are the escape zone: weak host immunity, like
a weak linezolid ceiling effect, is what lets resistant growth escape
control, so the switch-curve shading and histogram in step 13 use the same
low-is-bad convention as \texttt{MC\_emax\_lzd.py}, the reverse of
\texttt{MC\_ec50\_lzd.py}.

\subsection{Calibration history of \texttt{K\_IMMUNE\_MEAN} and \texttt{CV}}

The mean and $CV$ of the $k_{\mathrm{immune}}$ sampling distribution were
tuned, and later re-tuned, against the escape threshold rather than chosen
arbitrarily:

\begin{itemize}
\item At the model's own default, $k_{\mathrm{immune}} = 0.12$, only
$\sim$31--37\% of a sampled population resolved the infection.
\item The mean was first raised to $0.142$ so that $\sim$70\% of samples
would land above the escape threshold identified by
\texttt{ImmuneResponse\_threshold\_sweep.py}
($k_{\mathrm{immune}} = 0.1255$ h$^{-1}$) and resolve.
\item $CV$ was then widened from $0.20$ to $0.25$, which dropped resolution
to $\sim$64\%, so the mean was re-solved to $0.14719$ -- chosen so that the
30th percentile of $\mathrm{LogNormal}(\mathrm{mean}, CV{=}0.25)$ sits
exactly at the $0.1255$ threshold, restoring resolution to $\sim$70\%. This
value ($0.14719$) is what the module docstring at the top of the file still
documents.
\item Separately, $\rho_S$ was raised project-wide from $0.60$ to $0.61$
(see the $\rho_S$ sensitivity analysis referenced throughout the project).
Because of a vancomycin-driven competitive-release effect -- $S_b$ grows
larger before vancomycin is given, is then cleared by it, and leaves more of
blood's shared carrying capacity open for $R_b$ to claim -- this shifted the
$k_{\mathrm{immune}}$ escape threshold itself from $0.1255$ up to
$0.1521$ h$^{-1}$, a much larger sensitivity to the $\rho_S$ change than
either $EC_{50,L}$'s or $E_{max,L}$'s thresholds showed under the same
change.
\item That shift was not accompanied by a re-tuning of the sampling
distribution at the time: with the mean still at $0.14719$ but the threshold
now at $0.1521$, resolution dropped to only $\sim$40\% ($\sim$60\% escape).
\item \texttt{K\_IMMUNE\_MEAN} was therefore re-tuned a second time, to its
current code value of $0.172$ -- targeting $\sim$65\% of samples above the
(directly bisected, not estimated) $0.1521$ threshold, i.e.\ $\sim$35\%
relapse, within a stated 30--40\% target band.
\end{itemize}

\noindent \textbf{Note:} the module docstring at the top of
\texttt{MC\_immune\_response.py} still states
"$k_{\mathrm{immune}} \sim \mathrm{LogNormal}(\mathrm{mean}=0.14719,
CV=0.25)$," reflecting the \emph{first} re-tuning above. The code's actual
\texttt{K\_IMMUNE\_MEAN} is $0.172$, from the \emph{second} re-tuning
(triggered by the $\rho_S$ change) -- the docstring was not updated to match.
A reader relying on the docstring alone would compute the wrong
$\mu_{\ln}/\sigma_{\ln}$ for the distribution the code actually samples from.

\subsection{Why the other fixed parameters take their values}

\begin{itemize}
\item $\rho_R = 0.55$ is directly tuned to a $\sim$8.3\% fitness cost
relative to $\rho_S$, down from an earlier 20\% cost used elsewhere in the
project's history.
\item $\rho_{res,S} = 0.175$ was scaled $5\times$ (from an earlier $0.035$)
alongside the $\rho_S$ scaling.
\item $\rho_{res,R} = 0.1765$ is held in a narrow window just above the
reservoir's own persistence threshold ($0.17594055$, below which $R_{res}$
would clear on its own regardless of everything else) but below the point
at which $R_b$ would win the pretreatment competition for blood's shared
carrying capacity outright -- this is what allows $S_b$ to also establish a
visible blood infection in these simulations (see
\texttt{model\_Bacteremia.py}).
\item $B_{max,reservoir} = 10^4$ was lowered from an earlier $4.5\times10^6$
so that an escaped $R_{res}$ plateaus comfortably above the LOD but far
below its old, much larger equilibrium level.
\end{itemize}

\subsection{Why the switch-curve sweep range and re-derivation}

\texttt{SWEEP\_MIN}/\texttt{SWEEP\_MAX} $= [0.02, 0.22]$ is chosen to match
the range already swept by the standalone
\texttt{ImmuneResponse\_threshold\_sweep.py}, so the two scripts' derived
thresholds are directly comparable. \texttt{ESCAPE\_THRESH} is computed by
root-finding here rather than hardcoded, consistent with the same evolution
away from hardcoded threshold constants seen in
\texttt{MC\_ec50\_lzd.py} and \texttt{MC\_emax\_lzd.py} -- so that if
upstream parameters change again, the threshold used for shading and
population-splitting stays correct without manual updating.

\subsection{Why the population histogram is colored by simulated outcome}

The bottom panel of the switch-curve figure (step 13b) colors each sampled
patient by their own simulated final $R_{res}$ relative to the LOD, not by
which side of \texttt{ESCAPE\_THRESH} their sampled $k_{\mathrm{immune}}$
falls on. This mirrors the same design choice in \texttt{MC\_ec50\_lzd.py}
and \texttt{MC\_emax\_lzd.py}: it keeps the figure faithful to what the
Monte Carlo simulation actually produced, rather than to the idealized
deterministic switch, in case any samples land close enough to the threshold
that numerical or dynamical effects put them on the "wrong" side of it.

\end{document}
